\section{Population and Local Approximations}

To obtain a differentiable objective suitable for gradient-based optimization, we apply the standard Moreau-Yosida regularization of the constraint indicator, which for closed convex $\mathcal C$ yields the squared-distance penalty \citep{MoreauProx}:
\[
\min_\theta J(\theta)=
\min_\theta L_B(\theta) + \frac{\lambda}{2}\text{dist}(q(\theta), \mathcal{C})^2 = 
\min_\theta L_B(\theta) + \frac{\lambda}{2}\inf_{u \in \mathcal{C}} \ \text{dist}(q(\theta), u)^2.
\]
We call this the \textbf{Continual Learning as Optimization Constraint (CLOC)} formulation. GEM and A-GEM are similar in spirit in that they formulate preservation through constraints on old examples \citep{GEM,AGEM}. CLOC turns the old-task constraint into a proximal penalty built after task-$A$ training. In a population setting, e.g. when we have a dataset, we can translate this into a population risk objective by wrapping it in an expectation over the data and using sample-based analogues $\ell_B$ of $L_B$ and $q$ of $Q$, then taking the empirical average. Define the unscaled proximal loss as $p(a, \theta)=\text{dist}(q(a, \theta), \mathcal{C})^2$.
\[
\min_\theta
\;
\frac{1}{N_B}\sum_{i=1}^{N_B}\ell_B(b_i,\theta)
+
\frac{\lambda}{2N_A}\sum_{i=1}^{N_A}p(a_i, \theta)
\]
The two-distribution presentation is only for notational simplicity; additional old distributions can be handled by adding analogous constraints. We refer to this population version as \textbf{population CLOC (P-CLOC)}. One efficiency consideration is the penalty $p$: Ideally, we let the constraint contribute to each gradient step with its gradient. However, recomputing the full old-distribution penalty $p$ and its associated factors many times during optimization on distribution $B$ is usually prohibitively expensive, since it depends on the dataset $D_A$ associated with $A$. Additionally, reusing dataset $D_A$ later on does not respect the continual learning setting. For this reason, replace the exact penalty by a local $k$-th order Taylor-polynomial approximation around the old solution $\theta_A^\star$.
This turns the computation of $p$ into a repeated evaluation of precomputed derivative tensors and bounds memory requirements to $O(d^k)$ for a $d$-dimensional parameterization. Since memory scales exponentially with $k$, the standard scalable surrogate is a low-order Taylor approximation. We refer to this local approximation as \textbf{low-order CLOC (L-CLOC)}, and to its population version as \textbf{PL-CLOC}. It remains to specify the preserved quantity $q$ and constraint set $\mathcal{C}$. Notice that the gradient of $p$ is dependent on only one sample at once, which allows us to analyze $p$ and its gradients on a per-sample basis.
